\begin{center}
Problema C

Otimização de Rotas de Coleta de Solo - Simple Smart Agro

{\small Nome do arquivo: C.c, C.cpp, C.java, ou C.py}

{\large Timelimit: $2$}

{\tiny Por Alexandre da Silva}

\end{center}	

\footnotesize
			
A Simple Smart Agro está desenvolvendo um sistema automatizado para coleta de amostras de solo em uma fazenda. O sistema precisa otimizar as rotas de coleta para maximizar a eficiência do processo. Existem $N$ pontos específicos onde amostras de solo devem ser coletadas. O veículo de coleta parte sempre do depósito localizado na origem $(0, 0)$ e deve retornar ao mesmo ponto após coletar todas as amostras. Devido às características do terreno e dos equipamentos, o veículo só pode se mover nas direções cardeais (norte, sul, leste, oeste) e sempre em linha reta entre dois pontos consecutivos da rota. A distância entre dois pontos $(x_1, y_1)$ e $(x_2, y_2)$ é calculada usando a distância Manhattan: $|x_1 - x_2| + |y_1 - y_2|$.

\subsection*{Entrada}

\begin{itemize}
    \item A primeira linha contém um inteiro $N$ $(1 \leq N \leq 10)$, representando o número de pontos de coleta;
    \item As próximas $N$ linhas contêm dois inteiros $x_i$ e $y_i$ $(-1000 \leq x_i, y_i \leq 1000)$, representando as coordenadas do $i$-ésimo ponto de coleta;
       \item Todos os pontos de coleta são distintos entre si;
    \item Nenhum ponto de coleta coincide com a origem $(0,0)$.
\end{itemize}

\subsection*{Saída}

\begin{itemize}
    \item Uma única linha contendo um inteiro: a menor distância total (distância Manhattan) necessária para visitar todos os pontos de coleta e retornar à origem.
\end{itemize}

\subsection*{Exemplos}

\begin{center}
\begin{tabular}{|p{6cm}|p{4cm}|}
\hline
\textbf{Exemplo de Entrada} & \textbf{Exemplo de Saída} \\
\hline
\texttt{3} & \texttt{26} \\
\texttt{2 3} & \\
\texttt{7 1} & \\
\texttt{5 6} & \\
\hline
\texttt{4} & \texttt{42} \\
\texttt{3 4} & \\
\texttt{8 2} & \\
\texttt{12 8} & \\
\texttt{6 10} & \\
\hline
\texttt{1} & \texttt{8} \\
\texttt{2 2} & \\
\hline
\end{tabular}
\end{center}

\subsection*{Explicação do Primeiro Exemplo}

Partindo da origem $(0,0)$, temos os pontos de coleta em $(2,3)$, $(7,1)$ e $(5,6)$.
Uma rota ótima seria:
$(0,0) \rightarrow (2,3) \rightarrow (5,6) \rightarrow (7,1) \rightarrow (0,0)$

Calculando as distâncias:
\begin{itemize}
    \item $(0,0)$ para $(2,3)$: $|2-0| + |3-0| = 2 + 3 = 5$
    \item $(2,3)$ para $(5,6)$: $|5-2| + |6-3| = 3 + 3 = 6$
    \item $(5,6)$ para $(7,1)$: $|7-5| + |1-6| = 2 + 5 = 7$
    \item $(7,1)$ para $(0,0)$: $|0-7| + |0-1| = 7 + 1 = 8$
\end{itemize}

Distância total: $5 + 6 + 7 + 8 = 26$


\newpage

\begin{center}
\noindent
\lstinputlisting{abobora_25/C/C5.cpp}
\end{center}